\documentclass{beamer}
\usepackage[utf8]{inputenc}

\usetheme{Madrid}
\usecolortheme{default}
\usepackage{amsmath,amssymb,amsfonts,amsthm}
\usepackage{txfonts}
\usepackage{tkz-euclide}
\usepackage{listings}
\usepackage{adjustbox}
\usepackage{array}
\usepackage{tabularx}
\usepackage{gvv}
\usepackage{lmodern}
\usepackage{circuitikz}
\usepackage{tikz}
\usepackage{graphicx}

\setbeamertemplate{page number in head/foot}[totalframenumber]

\usepackage{tcolorbox}
\tcbuselibrary{minted,breakable,xparse,skins}



\definecolor{bg}{gray}{0.95}
\DeclareTCBListing{mintedbox}{O{}m!O{}}{%
	breakable=true,
	listing engine=minted,
	listing only,
	minted language=#2,
	minted style=default,
	minted options={%
		linenos,
		gobble=0,
		breaklines=true,
		breakafter=,,
		fontsize=\small,
		numbersep=8pt,
		#1},
	boxsep=0pt,
	left skip=0pt,
	right skip=0pt,
	left=25pt,
	right=0pt,
	top=3pt,
	bottom=3pt,
	arc=5pt,
	leftrule=0pt,
	rightrule=0pt,
	bottomrule=2pt,
	toprule=2pt,
	colback=bg,
	colframe=orange!70,
	enhanced,
	overlay={%
		\begin{tcbclipinterior}
			\fill[orange!20!white] (frame.south west) rectangle ([xshift=20pt]frame.north west);
	\end{tcbclipinterior}},
	#3,
}
\lstset{
	language=C,
	basicstyle=\ttfamily\small,
	keywordstyle=\color{blue},
	stringstyle=\color{orange},
	commentstyle=\color{green!60!black},
	numbers=left,
	numberstyle=\tiny\color{gray},
	breaklines=true,
	showstringspaces=false,
}
%------------------------------------------------------------
%This block of code defines the information to appear in the
%Title page
\title %optional
{12.164}
%\subtitle{A short story}

\author % (optional)
{RAVULA SHASHANK REDDY - EE25BTECH11047}

 \begin{document}
	
	
	\frame{\titlepage}
	\begin{frame}{Question}
    Let 
\begin{align*}
\vec{A} = \myvec{1 & 1 & 1 \\ 0 & 1 & 1 \\ 0 & 0 & 1}.
\end{align*}

Find the maximum number of linearly independent eigenvectors of \(\vec{A}\).
\end{frame}
\begin{frame}{Solution}
    Since $\vec{A}$ is a triangular matrix
\begin{align}
 & \lambda = 1,1,1 \\
(\vec{A} - \vec{I})\vec{x} &= \myvec{0 & 1 & 1 \\ 0 & 0 & 1 \\ 0 & 0 & 0} \vec{x} = 0, \quad \vec{x} = \myvec{x \\ y \\ z} \\
& y + z = 0 \implies y = -z \\
& z = 0 \implies y = 0 \\
\vec{x} &= \myvec{x \\ 0 \\ 0} = x \myvec{1 \\ 0 \\ 0}
\end{align}


Maximum number of linearly independent eigenvectors = 1


\end{frame}
\begin{frame}[fragile]
\frametitle{C Code}
\begin{lstlisting}
    #include <stdio.h>
#include <math.h>

int main() {
    // Matrix A
    double A[3][3] = {
        {1, 1, 1},
        {0, 1, 1},
        {0, 0, 1}
    };

    // Step 1: Since A is upper triangular, eigenvalues are diagonal elements
    double lambda[3];
    for (int i = 0; i < 3; i++) {
        lambda[i] = A[i][i];
    }

    printf("Eigenvalues:\n");
    \end{lstlisting}
    \end{frame}
\begin{frame}[fragile]
\frametitle{C Code}
\begin{lstlisting}
    for (int i = 0; i < 3; i++) {
        printf("λ%d = %.2f\n", i+1, lambda[i]);
    }

    // Step 2: All eigenvalues are equal (1), find eigenvectors for λ = 1
    // (A - I)
    double AI[3][3];
    for (int i = 0; i < 3; i++) {
        for (int j = 0; j < 3; j++) {
            AI[i][j] = A[i][j];
            if (i == j) AI[i][j] -= 1.0;
        }
    }

    printf("\nMatrix (A - I):\n");
    for (int i = 0; i < 3; i++) {
     \end{lstlisting}
    \end{frame}
\begin{frame}[fragile]
\frametitle{C Code}
\begin{lstlisting}
        for (int j = 0; j < 3; j++) {
            printf("%6.2f ", AI[i][j]);
        }
        
        printf("\n");
    }

    // From manual solving:
    // x3 = 0, x2 = 0, x1 = free → 1 eigenvector
    printf("\nEigenvector corresponding to λ = 1:\n");
    printf("v1 = [1 0 0]^T\n");

    printf("\nMaximum number of linearly independent eigenvectors = 1\n");
    return 0;
}

\end{lstlisting}
\end{frame}
\begin{frame}[fragile]
\frametitle{Python Direct}
\begin{lstlisting}
    import numpy as np

# Define matrix A
A = np.array([
    [1, 1, 1],
    [0, 1, 1],
    [0, 0, 1]
], dtype=float)

# Step 1: Eigenvalues
eigenvalues, _ = np.linalg.eig(A)
print("Eigenvalues:")
for i, val in enumerate(eigenvalues, start=1):
    print(f"λ{i} = {val:.2f}")

# Step 2: For eigenvalue = 1, compute (A - I)
I = np.eye(3)
AI = A - I
\end{lstlisting}
\end{frame}
\begin{frame}[fragile]
\frametitle{Python Direct}
\begin{lstlisting}
print("\nMatrix (A - I):")
print(AI)

# Step 3: Find null space to get eigenvectors
# (A - I)x = 0 → eigenvectors for λ = 1
u, s, vh = np.linalg.svd(AI)
tol = 1e-10
null_mask = (s <= tol)
null_space = np.compress(null_mask, vh, axis=0)
\end{lstlisting}
\end{frame}
\begin{frame}[fragile]
\frametitle{Python Direct}
\begin{lstlisting}
print("\nEigenvectors corresponding to λ = 1:")
if null_space.size == 0:
    print("No eigenvectors found.")
else:
    print(null_space.T)

# Step 4: Number of linearly independent eigenvectors
num_independent = null_space.shape[0]
print(f"\nMaximum number of linearly independent eigenvectors = {num_independent}")

\end{lstlisting}
\end{frame}

\end{document}
